\section{L8 / 8a -- GPU, CUDA \& OpenCL}

\subsection{GPU vs CPU \& the CUDA Model}

\mult{2}

\begin{concept}{GPU vs CPU}\\
    \begin{itemize}
        \item CPU: latency-optimised (few fast threads)
        \item GPU: \important{throughput} (many slow threads)
    \end{itemize}
\end{concept}

\begin{definition}{CUDA Hierarchy}\\
    \begin{itemize}
        \item kernel $\to$ thread (1/element)
        \item threads $\to$ block $\to$ grid
        \item block $\to$ \textbf{SM}; SM runs \textbf{warps of 32}
    \end{itemize}
\end{definition}

\multend

\begin{definition}{SM (Streaming Multiprocessor)}\\
    HW unit a block runs on. Fermi SM: 32 cores, 4 SFU, 16 L/S, 64KB shared/L1, warp scheduler. Grid of blocks distributed across SMs.
    % INSERT IMAGE: L8 slide 7 (grid/blocks) + slide 9 (Fermi SM)
\end{definition}

\begin{KR}{Recipe: Size a CUDA / OpenCL Launch}\\
    \textbf{One thread per element} \#threads = \#elements.
    \textbf{Block/grid} \texttt{threadsPerBlock} multiple of 32; grid $=\lceil$elements/block$\rceil$.
    \textbf{Index ranges} \texttt{[R][C]} $\to$ \texttt{threadIdx.x}$\in0..R{-}1$, \texttt{.y}$\in0..C{-}1$.
\end{KR}

\begin{example2}{Ex.\ S4--7 -- GPU vs CPU \& thread count}\\
    How do GPU/CPU differ? Threads for a $[10][5]$ op? threadID range?
    \tcblower
    CPU latency- vs GPU throughput-optimised. \important{$10\times5=50$ threads}; \texttt{threadIdx.x}$\in0..9$, \texttt{.y}$\in0..4$ (IDs $(0,0)..(9,4)$).
\end{example2}

\begin{example2}{Ex.\ S8--11 -- SM \& limitations}\\
    What is an SM? Why not one fine thread for everything?
    \tcblower
    SM = unit a block runs on (warps of 32). Limits: branch divergence serialises a warp, memory coalescing matters, too-fine work wastes scheduling.
\end{example2}

\subsection{Rendering Pipeline}

\begin{concept}{Stages}\\
    world coords $\to$ vertex shader $\to$ assembly $\to$ (geometry) $\to$ rasterise $\to$ fragment shader (per pixel) $\to$ blend. Mali: Utgard (tiling) $\to$ Midgard (OpenCL).
    % INSERT IMAGE: L8 slide 14 -- rendering pipeline
\end{concept}

\begin{example2}{Ex.\ S14/15 -- Why so parallel?}\\
    Why is rendering massively parallel?
    \tcblower
    Data \important{expands} from a few vertices to millions of pixels $\to$ per-pixel work is independent and huge.
\end{example2}

\subsection{OpenCL}

\mult{2}

\begin{definition}{Four Models}\\
    \begin{itemize}
        \item Platform (Host$\to$Device$\to$CU$\to$PE)
        \item Execution (work-item/group/NDRange)
        \item Memory (weak consistency)
        \item Programming
    \end{itemize}
\end{definition}

\begin{definition}{Memory Regions}\\
    \begin{itemize}
        \item Global, Constant
        \item Local (per work-group)
        \item Private (per work-item)
    \end{itemize}
\end{definition}

\multend

\begin{KR}{Recipe: Convert a Loop to a Kernel}\\
    drop the loop $\to$ \texttt{i=get\_global\_id(0)} $\to$ body on element \texttt{i}; work-items independent.
\end{KR}

\begin{examplecode}{OpenCL Vector Add}\\
\begin{lstlisting}[language=C, style=basesmol]
__kernel void vadd(__global float* a,
                   __global float* b,
                   __global float* c) {
    int i = get_global_id(0);
    c[i] = a[i] + b[i];   // one work-item per element
}
\end{lstlisting}
\end{examplecode}

\begin{example2}{Ex.\ S27/28 -- Loop $\to$ kernel}\\
    Rewrite \texttt{for(i<4096) c[i]=a[i]+b[i];} and give the NDRange.
    \tcblower
    See code box. NDRange = 4096 work-items (e.g.\ group 512); the loop is replaced by the index space.
\end{example2}

\begin{example2}{Ex.\ S21--26/31--33 -- Models \& sync}\\
    Name the four models; where can you synchronise?
    \tcblower
    Platform/Execution/Memory/Programming. Sync \important{only within a work-group} (barrier, fence, atomics); host enqueues in order, completion in/out-of-order.
\end{example2}

\subsection{GPU Scheduling}

\begin{concept}{Coarse vs Fine}\\
    \begin{itemize}
        \item coarse: blocks $\to$ SMs
        \item fine: warp scheduler interleaves 32-thread warps to \important{hide latency}
    \end{itemize}
\end{concept}

\begin{example2}{Ex.\ S29 -- Why interleave warps?}\\
    Why does the warp scheduler interleave?
    \tcblower
    When one warp stalls on memory, another ready warp runs $\to$ execution units stay busy (latency hiding).
\end{example2}

\subsection{AMP Recap}

\begin{example2}{Ex.\ S40/41 -- Different ISA under one OS?}\\
    \tcblower
    \important{No} -- separate OSes coupled by openAMP (\texttt{remoteproc} + \texttt{rpmsg}/\texttt{virtio}). Same as L7.
\end{example2}
